\section{Full proof of Theorem X.4.2 ($c_K$ leading + subleading identity)}

This appendix gives the full proof of Theorem X.4.2 of the main text:
under the hypotheses of Theorem X.4.1 (i.e., \(\chi\) primitive
non-principal of conductor \(q\), \(\rho = \tfrac12 + i\tau\) a \emph{simple}
zero of \(L(s,\chi)\) with \(\tau \ne 0\)), and assuming every
off-target nontrivial zero \(\rho' \ne \rho\) of \(L(s,\chi)\) in the
truncation height \(|\mathrm{Im}(\rho')| \le T(K)\) of the Inoue~\cite{Inoue2021}
explicit-formula box is also \emph{simple} (we write \(T(K)\) for the
truncation height, distinct from the partial prime-power sum
\(T_K(\chi,\rho)\) of \S\,X.1 of the main text; the \(o(1)\) rate further
uses RH for \(L(s,\chi)\), see \S\,B.3 and \S\,B.4),
\begin{equation}
\label{eq:cK-appendix}
c_K(\chi,\rho) \;=\; \frac{\log K}{L'(\rho,\chi)} \;+\; C_1(\chi,\rho) \;+\; o(1)
\qquad (K \to \infty),
\tag{$\dagger$}
\end{equation}
where \(C_1(\chi,\rho) = -L''(\rho,\chi) / (2\,L'(\rho,\chi)^2)\) and
\(c_K(\chi,\rho) = \sum_{n \le K} \mu(n)\,\chi(n)\,n^{-\rho}\).

The identity (\(\dagger\)) is \textbf{unconditional in \(\rho\)} given
simplicity of \(\rho\) and of the crossed off-target zeros. The \(o(1)\) rate
\(O(K^{-1/2 + \varepsilon})\) for every \(\varepsilon > 0\), and the
sharper rate
\(O\bigl(K^{-1/2} \exp\bigl((\log K)^{1/2}(\log\log K)^{14}\bigr)\bigr)\)
obtained via the character analogue of Soundararajan~\cite[Theorem~1]{Soundararajan2009},
are both conditional on RH for \(L(s, \chi)\). For the four characters
\(\chi_{-4}, \chi_5, \chi_{11}\) of \S\,X.5, the nontrivial zeros of
\(L(s,\chi)\) in the relevant explicit-formula range are numerically
verified to lie on the critical line (provenance in the Supplementary
computation audit), so these RH-conditional rates are the operative
ones for the finite \(K\) reported here; we do not assert them as
unconditional theorems. The unconditional fallback is the weaker
Akatsuka~\cite{Akatsuka2017EulerProduct} eq. (2.5) \(O(1/\log K)\) bound.

\subsection{Setup --- Perron formula for \(c_K\) via \(1/L(s, \chi)\)}

\subsubsection{The Dirichlet series for \(1/L(s, \chi)\)}

For \(\chi\) primitive of conductor \(q \ge 2\) and \(\mathrm{Re}(s) > 1\),
\[
\frac{1}{L(s, \chi)}
\;=\;
\sum_{n=1}^{\infty} \frac{\mu(n)\,\chi(n)}{n^s},
\]

absolutely convergent on \(\mathrm{Re}(s) > 1\). The Dirichlet
coefficients are \(a_n = \mu(n)\,\chi(n)\), and the partial sum
\(c_K(\chi, \rho) = \sum_{n \le K} a_n n^{-\rho}\) is the spectroscope
object (argument order matching \S\,X.1 of the main text).

\subsubsection{Perron formula (truncated)}

The standard truncated Perron formula (Tenenbaum, \emph{Introduction to
Analytic and Probabilistic Number Theory}, Theorem II.2.3; Inoue
(2021), \emph{Some explicit formulas for partial sums of Möbius
functions}, Journal de Théorie des Nombres de Bordeaux \textbf{33} (2021),
273--315, eq. (4.1) and Theorem 1) gives the following representation.

For any \(\sigma_0' > \tfrac12\) and \(T \ge 2\), with the convention that
\(c_K\) uses the half-weight at integer \(K\),
\begin{equation}
\label{eq:Perron-truncated}
c_K(\chi, \rho)
\;=\;
\frac{1}{2\pi i} \int_{\sigma_0' - iT}^{\sigma_0' + iT}
\frac{K^w}{w} \cdot \frac{1}{L(w + \rho, \chi)} \,dw
\;+\; \mathcal{R}'(K, T, \rho),
\end{equation}
where the kernel is \(K^w / w\) and the Dirichlet series factor is
\(1/L(w + \rho, \chi)\), obtained by the substitution \(s = w + \rho\).
The truncation error \(\mathcal{R}'(K, T, \rho)\) is bounded by Inoue
(2021) Theorem 1 via the \(J_1, J_2, J_3\) estimates of that paper:
\[
\mathcal{R}'(K, T, \rho)
\;\ll\;
\frac{K^{\sigma_0' - 1/2}\log K}{T}
\;+\;
\min\!\Bigl(1,\, \frac{K^{1/2}}{T\langle K\rangle}\Bigr),
\]

with \(\langle K \rangle\) the distance from \(K\) to the nearest integer.

\subsubsection{Inoue~\cite[Theorem~1]{Inoue2021} --- truncated explicit formula}

The following statement is verbatim from Inoue~\cite{Inoue2021}, arXiv:1805.05015v1, page 3:

\begin{quote}
\textbf{Theorem 1} (\cite{Inoue2021}). Let \(x > 0\), \(q \ge 2\),
\(T \ge \max\{T_0, \exp(q^{1/3}), 2/x\}\). Then, uniformly for all
primitive Dirichlet characters \(\chi\) modulo \(d\) with \(d \le q\),
there exists \(T_\nu \in [T, 2T]\) satisfying

\[M^*(x, \chi)
  \;=\; \sum_{|\gamma| < T_\nu} \frac{1}{(m(\rho)-1)!}\,
       \lim_{s \to \rho} \frac{d^{m(\rho)-1}}{ds^{m(\rho)-1}}
       \Bigl((s - \rho)^{m(\rho)}\,\frac{x^s}{L(s, \chi)\,s}\Bigr)
       \;+\; \mathop{\mathrm{Res}}_{s=0}\Bigl(\frac{x^s}{L(s, \chi)\,s}\Bigr)
       \;+\; \cdots\]
\end{quote}

This is the unshifted form, applied to the partial-character-Möbius
sum \(M^*(x, \chi) = \sum_{n \le x}' \mu(n)\,\chi(n)\). Dividing by
\(K^\rho\) and specializing at \(x = K\), \(m(\rho) = 1\) (the simplicity
hypothesis on \(\rho\)), and after the change of variable \(w = s - \rho\),
one obtains (\ref{eq:Perron-truncated}).

\subsection{Pole-structure analysis at \(w = 0\)}

The pole-structure computation of this section is independently
formalised in Lean 4 / Mathlib v4.28.0 as \texttt{local\_perron\_residue} in
\texttt{formal-conjectures/LocalPerronResidue.lean}. The Lean theorem
states the residue identity as a \texttt{Tendsto} limit at the canonical
form \(L\) analytic with a simple zero at \(0\) (the general-\(\rho\) case
of this appendix reduces to the canonical form by the substitution
\(L \mapsto L(\,\cdot\, + \rho)\)); the proof is \textbf{unconditional}
(0 \texttt{sorry}, only \texttt{propext}, \texttt{Classical.choice}, \texttt{Quot.sound}
axioms).

We work in the shifted variable \(w = s - \rho\). Define
\[
F(w) \;:=\; \frac{K^w}{w \cdot L(w + \rho, \chi)}.
\]

Since \(L(w + \rho, \chi)\) has a \emph{simple} zero at \(w = 0\) (by hypothesis),
\(1/L(w+\rho, \chi)\) has a simple pole at \(w = 0\). Combined with the
simple pole of \(K^w / w\) at \(w = 0\), the integrand \(F\) has a
\textbf{double pole at \(w = 0\)}.

\subsubsection{Laurent expansion of \(1/L(w+\rho, \chi)\)}

By Taylor expansion at the simple zero,
\[
L(w + \rho, \chi)
\;=\;
a_1 w + a_2 w^2 + a_3 w^3 + \cdots,
\qquad a_k \;=\; \frac{L^{(k)}(\rho, \chi)}{k!},\quad a_1 = L'(\rho, \chi) \ne 0.
\]

Therefore
\[
\frac{1}{L(w + \rho, \chi)}
\;=\;
\frac{1}{a_1 w} \cdot \frac{1}{1 + (a_2/a_1) w + (a_3/a_1) w^2 + \cdots}.
\]

Expanding the geometric factor as \((1 + u)^{-1} = 1 - u + u^2 - \cdots\)
and substituting \(a_1 = L'(\rho, \chi)\), \(a_2 = L''(\rho, \chi)/2\):
\begin{equation}
\label{eq:invL-Laurent}
\frac{1}{L(w + \rho, \chi)}
\;=\;
\frac{1}{L'(\rho, \chi)\,w}
\;-\;
\frac{L''(\rho, \chi)}{2\,L'(\rho, \chi)^2}
\;+\;
O(w).
\end{equation}

\subsubsection{Laurent expansion of \(K^w / w\)}

Trivially,
\[
\frac{K^w}{w}
\;=\;
\frac{1}{w}\bigl(1 + (\log K)\,w + \tfrac{1}{2}(\log K)^2 w^2 + O(w^3)\bigr)
\;=\;
\frac{1}{w} \;+\; \log K \;+\; \tfrac{1}{2}(\log K)^2 \,w \;+\; O(w^2).
\]

\subsubsection{Product expansion and the double-pole residue}

Multiplying the two expansions:
\[
F(w)
\;=\;
\Bigl[\tfrac{1}{w} + \log K + \tfrac{1}{2}(\log K)^2 w + \cdots\Bigr]
\cdot
\Bigl[\tfrac{1}{L'(\rho,\chi)\,w} - \tfrac{L''(\rho,\chi)}{2 L'(\rho,\chi)^2} + O(w)\Bigr].
\]

Collecting powers of \(w\):
- \(w^{-2}\) coefficient: \(\tfrac{1}{L'(\rho,\chi)}\) (from \(\tfrac{1}{w} \cdot \tfrac{1}{L'(\rho,\chi)\,w}\));
- \(w^{-1}\) coefficient: \(\tfrac{\log K}{L'(\rho,\chi)} - \tfrac{L''(\rho,\chi)}{2\,L'(\rho,\chi)^2}\) (cross-terms).

Hence
\begin{equation}
\label{eq:double-pole-residue}
\mathop{\mathrm{Res}}_{w = 0} F(w)
\;=\;
\frac{\log K}{L'(\rho,\chi)}
\;-\;
\frac{L''(\rho, \chi)}{2\,L'(\rho, \chi)^2}
\;=\;
\frac{\log K}{L'(\rho,\chi)} + C_1(\chi, \rho).
\end{equation}

This is \textbf{Lemma X.3.1 of the main text}, established algebraically;
the local part of (\ref{eq:cK-appendix}) is now done. It remains to
control the off-target zero residues, the trivial-zero residues, and
the contour pieces.

\subsection{Off-target zero residues, trivial-zero residues, and contour pieces}

We now shift the Perron contour from \(\mathrm{Re}(w) = \sigma_0' > \tfrac12\)
leftward to a zero-avoiding line \(\mathrm{Re}(w) = -A\) for some \(A > 0\),
crossing all the nontrivial zeros of \(L(s, \chi)\) in the strip
\(-A \le \mathrm{Re}(w) \le \sigma_0'\). The integrand \(F(w)\) has poles
at \(w = 0\) (the target double pole, with residue
(\ref{eq:double-pole-residue})), at \(w = \rho' - \rho\) for every
other nontrivial zero \(\rho' \ne \rho\) of \(L(s, \chi)\), at trivial
zeros, and at \(w = -\rho\) from a (potential) pole of \(L(s, \chi)\) at
\(s = 0\) --- but \(L(s, \chi)\) is entire for primitive non-principal
\(\chi\), so the pole at \(w = -\rho\) from \(1/L\) is absent. We do pick
up the term \(\mathop{\mathrm{Res}}_{s = 0}\bigl(x^s/(L(s,\chi)\,s)\bigr)\) from Inoue's formula
(this is the \(s = 0\) residue in the unshifted form; after dividing by
\(K^\rho\) it contributes a term of magnitude \(K^{-\rho}/(L(0,\chi)) =
O(K^{-1/2})\) in modulus, \(o(1)\)).

\subsubsection{Off-target nontrivial-zero residues}

For each \(\rho' \ne \rho\) a nontrivial zero of \(L(s, \chi)\), write
\(w_{\rho'} = \rho' - \rho\). Assuming \(\rho'\) is \emph{simple} (which is
generically expected and is the regime relevant for the leading +
subleading identity), the residue of \(F\) at \(w = w_{\rho'}\) is
\[
\mathop{\mathrm{Res}}_{w = w_{\rho'}} F(w)
\;=\;
\frac{K^{w_{\rho'}}}{w_{\rho'} \cdot L'(\rho', \chi)}
\;=\;
\frac{K^{\rho' - \rho}}{(\rho' - \rho)\,L'(\rho', \chi)}.
\]

Under RH for \(L(s, \chi)\), \(|K^{\rho' - \rho}| = 1\) for all
\(\rho' = \tfrac12 + i\gamma'\) with \(\gamma' \in \mathbb{R}\). The
aggregate off-target sum (over zeros \(\rho'\) in the strip
\(|\gamma'| \le T\)) is then bounded by
\[
\Bigl|\!\sum_{\rho' \ne \rho,\ |\gamma'| \le T}
\frac{K^{i(\gamma' - \tau)}}{(\rho' - \rho)\,L'(\rho', \chi)}\Bigr|
\;\ll\;
\sum_{\rho' \ne \rho,\ |\gamma'| \le T}
\frac{1}{|\gamma' - \tau|\,|L'(\rho', \chi)|}.
\]

This is the off-target \emph{aggregate} whose control determines the
\(o(1)\) rate.

\subsubsection{Trivial-zero residues}

Trivial zeros of \(L(s, \chi)\) are at \(s = -2k\) (\(\chi\) even) or
\(s = -2k - 1\) (\(\chi\) odd), \(k \ge 0\). After the shift \(w = s - \rho\),
these contribute residues of size \(K^{-\mathrm{Re}(\rho) - 2k}
= K^{-1/2 - 2k}\), summable to \(O(K^{-1/2})\), which is \(o(1)\).

\subsubsection{Shifted contour pieces}

The shifted contour \(\mathrm{Re}(w) = -A\) contributes, on a suitable
zero-avoiding truncation height \(T(K) \in [K/(\log K)^B, 2K/(\log K)^B]\), an
integral bounded by (using a reciprocal-\(L\) bound on the vertical
shifted line)
\[
\Bigl|\frac{1}{2\pi i}\int_{(-A)}\!\frac{K^w}{w\,L(w+\rho,\chi)}\,dw\Bigr|
\;\ll\;
K^{-A} \cdot T(K) \cdot \exp\bigl(C\,(\log\log T(K))^2\bigr).
\]

For \(A > 1/2\) and \(T(K) \le K/(\log K)^B\) with \(B\) sufficiently large,
this is \(K^{-A + 1 + o(1)} \cdot \exp(O((\log\log K)^2))\), which is
\(o(1)\).

Horizontal pieces give factors like \((T(K)/K)^A \cdot \mathrm{polylog}\),
\(o(1)\) for \(T(K) \le K (\log K)^{-B}\) and large enough \(B\).

\subsubsection{The Perron truncation error}

By Inoue~\cite[Theorem~1]{Inoue2021}, the truncation error
\(\mathcal{R}'(K, T(K), \rho)\) is bounded by
\[
\mathcal{R}'(K, T(K), \rho)
\;\ll\;
\frac{K^{\sigma_0' - 1/2} \log K}{T(K)}
\;+\;
\min\!\Bigl(1,\,\frac{K^{1/2}}{T(K) \langle K\rangle}\Bigr).
\]

For \(T(K) \asymp K (\log K)^{-B}\) and \(\sigma_0' = 1/2 + 1/\log K\), both
pieces are \(O(K^{-1/2 + \varepsilon}) \cdot (\log K)^{O(1)}\) for any
\(\varepsilon > 0\), in particular \(o(1)\).

\subsubsection{Assembly}

Combining the local residue (\ref{eq:double-pole-residue}) with the
off-target and contour pieces:
\begin{multline*}
c_K(\chi, \rho)
\;=\;
\frac{\log K}{L'(\rho, \chi)}
\;+\;
C_1(\chi, \rho)
\;+\;
\bigl(\text{off-target aggregate}\bigr)
\;+\;
O(K^{-1/2})
\\
\;+\;
O\!\bigl(K^{-A+1+o(1)}\exp(O((\log\log K)^2))\bigr)
\;+\;
\mathcal{R}'(K, T(K), \rho).
\end{multline*}
All terms in the last three parentheses are \(o(1)\). The off-target
aggregate is also \(o(1)\) under RH for \(L(s, \chi)\), by partial
summation against Soundararajan~\cite{Soundararajan2009}, \emph{Partial sums of the Möbius
function}, J. reine angew. Math. (Crelle's Journal) \textbf{631} (2009),
141--152, Theorem 1:
\[
\text{(\cite{Soundararajan2009})}
\quad
\text{under RH:}\quad
M(x) \;\ll\; \sqrt x \,\exp\!\bigl((\log x)^{1/2}\,(\log\log x)^{14}\bigr).
\]

Translating to \(M^*(x, \chi)\) (the character analogue) gives the
required \(o(1)\) rate for the off-target aggregate under RH for
\(L(s, \chi)\).

This establishes (\(\dagger\)). \(\hfill\square\)

\subsection{Rate of the \(o(1)\)}

The proof above yields the following rates.

{\def\LTcaptype{none} % do not increment counter
\begin{longtable}[]{@{}
  >{\raggedright\arraybackslash}p{(\linewidth - 2\tabcolsep) * \real{0.5000}}
  >{\raggedright\arraybackslash}p{(\linewidth - 2\tabcolsep) * \real{0.5000}}@{}}
\toprule\noalign{}
\begin{minipage}[b]{\linewidth}\raggedright
Rate
\end{minipage} & \begin{minipage}[b]{\linewidth}\raggedright
Hypotheses
\end{minipage} \\
\midrule\noalign{}
\endhead
\bottomrule\noalign{}
\endlastfoot
\(o(1) = O\bigl(K^{-1/2 + \varepsilon}\bigr)\) for every \(\varepsilon > 0\) & RH for \(L(s, \chi)\). \\
\(o(1) = O\bigl(K^{-1/2}\exp\bigl((\log K)^{1/2}(\log\log K)^{14}\bigr)\bigr)\) & RH for \(L(s, \chi)\), via the character analogue of Soundararajan~\cite[Theorem~1]{Soundararajan2009}. The operative rate for \S\,X.5: for \(\chi_{-4}, \chi_5, \chi_{11}\) the nontrivial zeros of \(L(s,\chi)\) in the relevant explicit-formula range are numerically verified on the critical line (Supplementary computation audit); not asserted as an unconditional theorem. \\
\(o(1) = O\bigl(\log K / \sqrt K\bigr)\) & RH plus a Gonek--Hejhal-type bound of the form \(\sum_{0 < \gamma < T} |L'(\rho, \chi)|^{-2} \ll T (\log T)^{O(1)}\). \\
\end{longtable}
}

In the numerical work of \S\,X.5.3, the empirical residual
\(R(K) := c_K - \log K / L'(\rho, \chi) - C_1(\chi, \rho)\) at
\(K = 2\cdot 10^6\) falls in the interval \([0.027,\,0.375]\) across the
four pairs; this is consistent with the Soundararajan-conditional rate
within an \(O(1)\) implicit-constant factor (specifically,
\(|R(K)| / (\log K / \sqrt K) \in [3,\,36]\), well within the typical
\(\exp\bigl(C\,(\log K)^{1/2}(\log\log K)^{14}\bigr)\) envelope at this
scale).

\subsection{What the proof does \emph{not} claim}

For clarity:

\begin{itemize}
\item
  The proof does \textbf{not} establish the Aoki--Koyama--Mertens limit
  for \(E_K \log K\), which is the \emph{multiplicative-side} asymptotic.
\item
  The proof does \textbf{not} establish the \emph{full} Perron-leading theorem
  \(c_K = \log K / L'(\rho) + o(\log K)\) unconditionally; the
  off-target aggregate's \(o(\log K)\) control (which is \emph{weaker} than
  the \(o(1)\) obtained in (\(\dagger\)) at the \emph{leading + subleading}
  level) depends on RH for \(L(s, \chi)\) and is what is required for
  composing with (AK) to obtain (NDC).
\item
  The proof does \textbf{not} address the possibility of off-target
  \emph{multiple} zeros, which under DRH/EDRH alone are not excluded;
  the manuscript's open Question Q:Perron asks for either
  the off-target multiple-zero residue control or a sufficient
  ``all crossed off-target zeros simple'' hypothesis.
\end{itemize}
